\textbf{ImageNet-200．}構成を固定したまま単体で評価した（\NInTwoZeroZeroNDraws{}通りの抽出の平均，較正画像$n{=}800$；付録\ref{sec:app_in200}の\tabref{in200}）．$\sigma$はnear-OODで\AUC{} \NInTwoZeroZeroMainNearAuroc\%，\FPR{} \NInTwoZeroZeroMainNearFpr\%，far-OODで\NInTwoZeroZeroMainFarAuroc\%，\NInTwoZeroZeroMainFarFpr\%であり，OpenOOD v1.5~\cite{zhang2023openood}の報告値の最良（near \NInTwoZeroZeroBestPubNearAuroc\%／\NInTwoZeroZeroBestPubNearFpr\%，far \NInTwoZeroZeroBestPubFarAuroc\%／\NInTwoZeroZeroBestPubFarFpr\%；ImageNet-200で学習したResNet-18）を大きく上回る．ラベルの数を揃えない基線として，同じ2つの特徴で全学習画像258,951枚を使うkNN（$k{=}50$）とクラス平均，CLIP ViT-B/16のゼロショットのMCMとNegLabelも実測した．kNNは\FPR{} \NInTwoZeroZeroKnnAllKFiveZeroNearFpr\%／\NInTwoZeroZeroKnnAllKFiveZeroFarFpr\%，クラス平均は\NInTwoZeroZeroNcmAllNearFpr\%／\NInTwoZeroZeroNcmAllFarFpr\%で16枚の静的$p$（\NInTwoZeroZeroStaticNearFpr\%／\NInTwoZeroZeroStaticFarFpr\%）と同程度にとどまり，MCMとNegLabelは\NInTwoZeroZeroMcmClipNearFpr\%／\NInTwoZeroZeroMcmClipFarFpr\%，\NInTwoZeroZeroNeglabelClipNearFpr\%／\NInTwoZeroZeroNeglabelClipFarFpr\%である．ラベル付き画像を約80倍使っても伝播の分（$\pP$だけで\NInTwoZeroZeroLpNearFpr\%／\NInTwoZeroZeroLpFarFpr\%）は得られない．一方，上限なしの種は効かない．LGC-maxは\NInTwoZeroZeroLgcNearFpr\%／\NInTwoZeroZeroLgcFarFpr\%で，$\pN$を加えた本稿の構成はこれらより悪い．IDが9,000枚しかなく，SSB-hardでは未知が全画像の\NInTwoZeroZeroMixSsbHard\%を占めて折の全画像が種になる（付録\ref{sec:app_in200}；単純BH法の種では\NInTwoZeroZeroBhNearFpr\%）．上限付きの種（\secref{cap}）では\NInTwoZeroZeroStoreycNearFpr\%／\NInTwoZeroZeroStoreycFarFpr\%と$\pP$だけと同等以上になる．

\textbf{CIFAR-100．}OpenOOD v1.5のCIFAR-100ベンチマーク（\secref{exp}）で，同じ凍結特徴と学習画像からの16枚（クラスあたり，重ならない6通りの抽出，較正画像$n{=}400$），同じ構成で単体の値を測った（\tabref{cifar}）．$\sigma$はnear-OODで\AUC{} \NCifarMainNearAuroc\%，\FPR{} \NCifarMainNearFpr\%，far-OODで\NCifarMainFarAuroc\%，\NCifarMainFarFpr\%で，報告値（CIFAR-100で学習したResNet-18）の最良であるOE（near \NCifarBestPubNearAuroc\%／\NCifarBestPubNearFpr\%，外れ値データで学習）とRotPred（far \NCifarBestPubFarAuroc\%／\NCifarBestPubFarFpr\%）を\AUC{}で10ポイント以上上回る．全学習画像50,000枚のkNN（$k{=}50$，\NCifarKnnAllKFiveZeroNearFpr\%／\NCifarKnnAllKFiveZeroFarFpr\%）とクラス平均（\NCifarNcmAllNearFpr\%／\NCifarNcmAllFarFpr\%）は16枚の静的$p$（\NCifarStaticNearFpr\%／\NCifarStaticFarFpr\%）から数ポイント以内にとどまり，$32\times32$の画像ではゼロショットのMCMとNegLabelは弱い（\NCifarMcmClipNearFpr\%／\NCifarMcmClipFarFpr\%，\NCifarNeglabelClipNearFpr\%／\NCifarNeglabelClipFarFpr\%）．ここでも利得は$\pP$の伝播にあり，種は効かない．$\pP$だけで\NCifarLpNearFpr\%／\NCifarLpFarFpr\%，$\sigma$との差はnear-OODで\NCifarMainNearDfprVsLp{} [\NCifarMainNearDfprVsLpLo, +\NCifarMainNearDfprVsLpHi]，far-OODで+\NCifarMainFarDfprVsLp{} [+\NCifarMainFarDfprVsLpLo, +\NCifarMainFarDfprVsLpHi]ポイントで，far-OODでは$\pN$が害になる．未知がMNIST・SVHN・Places365で全画像の\NCifarMixFarHiMin{}--\NCifarMixFarHiMax\%を占め，割合$q$の種がIDのほぼ全数になって$\pN$は情報を失う（付録\ref{sec:app_in200}）．上限付きの種では近\NCifarStoreycNearFpr\%，遠\NCifarStoreycFarFpr\%になる（\secref{cap}）．
